\PID
Отпорност хигристора, нелинеарног сензора влажности, позната је таблично за разне вредности влажности и температуре, као што је 
приказано у Табели \ref{tab:\ID:in}. Због сложеног понашања овог сензора, у мерном систему у коме се он користи користи се још
додатно и сензор температуре. Сензор температуре је измерио вредност $t = 33\rm^\circ C$. Израчунати влажност коју показује хигристор
ако је његова измерена отпорност $R = 100\unit{k\Omega}$.

\begin{table}
\centering  
\begin{tabular}{c|cccc}
\hline
${\rm RH}\unit{[\%]}\backslash t\unit{[^\circ C]}$ & 20 & 30 & 40 & 50 \\
\hline
20\% & 6,7 $\unit{M\Omega}$ & 3.9 $\unit{M\Omega}$ & 2,4 $\unit{M\Omega}$ & 1,45 $\unit{M\Omega}$ \\
40\% & 420 $\unit{k\Omega}$ & 235 $\unit{k\Omega}$ & 140 $\unit{k\Omega}$ & 88 $\unit{k\Omega}$ \\
60\% & 39 $\unit{k\Omega}$ & 25 $\unit{k\Omega}$ & 17,5 $\unit{k\Omega}$ & 13 $\unit{k\Omega}$ \\
80\% & 7,0 $\unit{k\Omega}$ & 5,0 $\unit{k\Omega}$ & 3,9 $\unit{k\Omega}$ & 3,0 $\unit{k\Omega}$ \\
\hline
\end{tabular}
\caption{Отпорнсот хигристора у задатим условима.}
\label{tab:\ID:in}
\end{table}